\section{The search tree in Rust}
\label{sec:tree}

\subsection{The representational problem}

A search tree is a recursively owned structure: a node owns its children, and
each child owns its own children. In a language with a tracing garbage collector
this is unremarkable. In Rust it collides directly with the ownership model,
because the shape that is natural to draw is the shape that is hardest to express:
a mutable reference from a node to its child, while the parent is still borrowed, is
exactly what the borrow checker forbids.

The standard solutions, and why the choice matters, are laid out in
Table~\ref{tab:representations}. The important observation is that the tree is
mutated from the root downwards during selection and from the leaves upwards
during backup, that it is traversed millions of times over a training run, and
that its nodes are small. These three facts point at one representation.

\begin{table}[t]
\centering
\small
\caption{Tree representations in Rust. The arena is the standard choice for
game-tree engines for the reasons in the right-hand column.}
\label{tab:representations}
\begin{tabularx}{\textwidth}{@{}l X@{}}
\toprule
Representation & Consequences \\
\midrule
\code{Box<Node>} children & Recursive ownership. Descent requires holding a
\code{\&mut} to the parent while borrowing the child: not expressible. Descent
becomes a loop of \code{take}/\code{replace} tricks, and each node is a separate
allocation with pointer-chasing access. \\
\code{Rc<RefCell<Node>{}>} & Compiles and is pleasant to read, at a price: runtime
borrow checks on every access, reference-count traffic on every clone, and a
\code{RefCell} borrow held across a search step can panic. Debugging is harder
because a violation is a runtime panic instead of a compile error. \\
Raw pointers with \code{unsafe} & No borrow checks, no reference counting,
fastest in principle. The safety argument is entirely manual and unverifiable, in
a component that runs a few thousand times per move. \\
Flat arena
(one \code{Vec}, integer indices) & One allocation region for the whole tree,
cache-friendly traversal, no reference counting, no borrow conflicts, and
trivially serializable and inspectable. The cost is that the compiler no longer
proves that an index is valid, and that dangling indices are possible in
principle. Both are manageable because the tree is append-only within one search
(Section~\ref{sec:invariants}). \\
\bottomrule
\end{tabularx}
\end{table}

\subsection{The arena design}
\label{sec:arena}

The tree is two vectors and two integers: a vector of nodes, a vector of edges
shared by all nodes, a root index, and the number of nodes. Edges refer to
children by index, not by pointer.

\begin{lstlisting}[style=rust,caption={The arena. Node ids and edge ids are
integers; no node owns another.},label={lst:arena}]
pub type NodeId = u32;
pub type EdgeId = u32;

pub struct Tree {
    nodes: Vec<Node>,
    edges: Vec<Edge>,     // all nodes' edges in one region
    root: NodeId,
    gen: u32,             // the network generation that produced the priors
}

pub struct Node {
    first_edge: EdgeId,   // contiguous run of this node's edges
    edge_count: u16,      // at most 225 legal moves fit in u16
    parent: NodeId,       // unused at the root; kept for diagnostics only
    visits: u32,          // sum of the children's visit counts
    terminal: Option<i8>, // Some(z) if the position is over, in [-1, 1]
}

pub struct Edge {
    mv: Move,             // 0..225, one byte
    prior: f32,           // the network's P(s, a); never changes
    visits: u32,          // N(s, a)
    value_sum: f32,       // W(s, a), accumulated from the chooser's view
    child: NodeId,        // meaning depends on the current generation
}
\end{lstlisting}

Two details of this layout carry more weight than they appear to.

\paragraph{Edges are contiguous and indexed by a range.} A node stores
\code{first\_edge} and \code{edge\_count} rather than a \code{Vec<Edge>}. Child
ranges of a node created in one expansion are appended together as a single run.
Expansion therefore performs one growth of one vector for an entire node, rather
than one allocation per node -- and it also means that a node's edges are
adjacent in memory, which is where the traversal time is spent. The edge count
fits in \code{u16} because the board has 225 cells; the size of a node is
therefore fixed and small.

\paragraph{\code{child} means ``child, if this node has been expanded''.} A
freshly created edge has no child node yet: the node beyond it has never been
visited, so it has never been evaluated, so there is nothing to store. The usual
Rust encodings are \code{Option<NodeId>} or a sentinel. Both are workable; the
sentinel is faster and easy to get wrong. If a sentinel is used, make it a named
constant that no valid index can equal, and assert against it in debug builds.

The generation field deserves a word. A prior belongs to the network that
produced it. During self-play the network is replaced between iterations, so a
tree that survives a network update -- or a tree reused across moves after an
update -- can mix priors from two different networks. Storing the generation with
the tree and discarding trees from older generations removes the possibility
silently. This is a correctness invariant that costs one integer.

\subsection{Invariants}
\label{sec:invariants}

An arena trades compile-time safety for a set of invariants that must be
maintained. They should be written down, because they are the only reason an
integer index is as safe as a reference. For a single search:

\begin{enumerate}
  \item \textbf{Ids are valid.} Every \code{NodeId} and \code{EdgeId} stored in the
        arena is in range for that arena. This holds because both vectors only
        grow, and no element is ever removed.
  \item \textbf{The structure is a tree rooted at \code{root}.} Every node except
        the root is reachable from the root through exactly one path of edges.
        The board can be reconstructed from that path, which is how
        Section~\ref{sec:walk} avoids storing a board per node.
  \item \textbf{Edges are append-only and children are allocated once.} An edge's
        child index, once set, is never changed within a generation. Repeated
        visits to the same node reuse the existing child, which is what makes the
        tree a tree and not a graph of duplicated subtrees.
  \item \textbf{Priors are normalized over the legal moves of the node} and sum to
        one (up to floating-point error) before root noise is applied.
  \item \textbf{A terminal node has no edges.} Terminality is decided by the
        rules engine before expansion, so a terminal node never acquires children
        and never receives a network evaluation.
\end{enumerate}

Invariant 1 is the interesting one, because it is where an arena is weaker than
\code{Rc}. It rules out \emph{generational indices} -- the standard technique for
detecting stale references in an arena that \emph{does} free slots. No slot is
freed here, so there is nothing to detect, and adding a generation counter to
every id would only widen the id and slow the traversal. This is a case where the
absence of deletion is what makes the representation simple: the invariant is
maintained by construction rather than by checking.

Invariant 3 is the one to guard, because it is where a plausible optimization
goes wrong. Merging two edges that lead to the same position -- transposition
detection, using the engine's incremental position hash -- would turn the tree
into a directed acyclic graph. That is a real technique in classical search, and
Section~\ref{sec:excluded} explains why it is deliberately not used here.

\subsection{Descending: the walk, and what is borrowed}
\label{sec:walk}

Selection descends from the root by repeatedly choosing an edge and following its
child. Three things change as the walk proceeds -- the current node, the board,
and the path used for backup -- and the borrow discipline follows from keeping
them separate.

\begin{lstlisting}[style=rust,caption={Selection. The borrow of the node ends
before the next iteration, so no conflict is possible.},label={lst:select}]
fn select(&self, board: &mut Board, path: &mut Vec<EdgeId>) -> NodeId {
    let mut node_id = self.root;
    loop {
        let node = &self.nodes[node_id as usize];   // immutable borrow ...
        if node.terminal.is_some() || node.edge_count == 0 {
            return node_id;                          // leaf: terminal or unexpanded
        }
        let edge_id = self.puct_argmax(node);        // reads only
        let edge = &self.edges[edge_id as usize];
        path.push(edge_id);
        board.play(edge.mv);                         // board is a separate owner
        node_id = edge.child;                        // ... borrow ends here
    }
}
\end{lstlisting}

Two properties make this loop compile and stay fast. The borrow of
\code{self.nodes[node\_id]} ends at the last use, which is before the next
iteration, because nothing derived from it is held across the loop body. And the
board is not stored in the tree at all: it is passed in as a separate
\code{\&mut Board}, mutated on the way down and unwound on the way up. The tree
owns moves and statistics; the board owns stones.

This separation is worth defending, because the alternative -- caching a board or
a hash in each node -- is tempting for debugging and costs memory and copy time
on paths that are already the hot loop.

\paragraph{Clone or undo during the walk?} Both are available: the rules engine
can clone a board cheaply (a few machine words, since the position is a set of
bitboards), or it can play and undo moves along the path. Three considerations
decide it. Cloning in place gives the evaluator an owned board and removes a
class of lifetime problems; \code{play}/\code{undo} avoids the copy but requires
the walk to be mirrored on the way back up, which is exactly where sign and
ordering mistakes happen. For a board that is a handful of \code{u64} words, the
clone is a few nanoseconds, which is small next to a network evaluation but not
small next to the PUCT arithmetic. The measurement that settles it is a
profile, not an argument; both are legitimate, and the choice should be
revisited only with a profile in hand.

\subsection{Expansion and the normalization order}
\label{sec:expand}

Expansion happens once per leaf, and it has four steps in a fixed order.

\begin{enumerate}
  \item \textbf{Terminal check first.} Ask the rules engine whether the position
        is over. If it is, store the exact outcome as the node's value and stop.
        No network call, and no edges. This ordering is not an optimization
        detail; it is what makes exact values propagate (a search can learn that a
        move forces a win precisely because the terminal values along the winning
        line are exact).
  \item \textbf{Encode and evaluate the leaf.} The encoder turns the board into
        plain arrays; the evaluator returns a policy over all cells and a value.
  \item \textbf{Mask, then normalize.} Discard the policy entries for illegal
        cells, then normalize the rest with a softmax.
  \item \textbf{Create the edges.} One edge per legal move, carrying the
        normalized prior and zeroed statistics.
\end{enumerate}

Step 3 is the classic defect site, and the failure mode is quiet rather than
loud. If the softmax is applied before masking, the probability mass of the
illegal cells is not removed, it is redistributed over the legal cells by the
renormalization --- and the resulting priors still sum to one, so nothing looks
wrong. The search is then biased in a way that no test of the sum will catch. The
rule is worth stating as a rule: \emph{the normalization is over the legal move
set, not over the network's output}.

The same care applies to the value. The network's $v$ is quoted from the
perspective of the side to move at the leaf; the edge statistics are quoted from
the perspective of the player who chose at the parent. One of the two must be
negated, and the choice must be made once and written down (Section~\ref{sec:backup}).

\paragraph{Allocation.} One growth of the edges vector for an entire node keeps
the allocation count proportional to the number of expanded nodes rather than to
the number of edges. Reserving capacity for the expected node count at the start
of a search removes even that. A search with a budget of a few hundred
simulations creates a few hundred to a few thousand nodes; the whole tree is a
few hundred kilobytes, and no allocation is on a critical path.

\subsection{Backup, and the sign convention}
\label{sec:backup}

Backup walks the recorded path in reverse and updates each edge:

\begin{lstlisting}[style=rust,caption={Backup. The value alternates sign at every
ply, because the players alternate.},label={lst:backup}]
fn backup(&mut self, path: &[EdgeId], leaf_value: f32) {
    let mut v = leaf_value;                 // from the view of the player to
    for &edge_id in path.iter().rev() {     // move AT THE LEAF
        let e = &mut self.edges[edge_id as usize];
        e.visits += 1;
        e.value_sum += v;
        v = -v;                             // the parent's chooser sees the opposite
    }
}
\end{lstlisting}

Two equivalent conventions exist, and mixing them is the defect. Either
\code{value\_sum} is accumulated from the perspective of the player who chose at
the parent node, in which case the value is negated once before the loop and
again at every step as above; or it is accumulated from the perspective of the
side to move at the child, in which case the negation happens implicitly when
$Q$ is read. Both are correct. A codebase that uses one of them in selection and
the other in backup is wrong, and the error is silent: the search still runs, it
merely prefers the wrong moves, more reliably against stronger opponents.

Three defences are worth building, and none of them is the type system. A
newtype wrapper for values does not catch a missing negation, because $v$ and
$-v$ have the same type; the type system can only catch the case where two
different \emph{logical} quantities share a representation, as in the
absolute-colour versus me/you distinction of Section~\ref{sec:solved}. What does
work is a hand-computed test: build a small tree by hand, set the leaf values,
call backup, and compare every $Q$ against a value computed on paper. A second
defence is the tactical suite of Section~\ref{sec:alphaepsilon}: a sign error in
backup makes a search fail to play a win in one, which is immediately visible and
is the reason the mock evaluator exists before the network does. A third is an
invariant that is free to assert: at any node, the weighted average of the
children's $Q$ values, weighted by their visit counts, must lie in $[-1,1]$,
because every value fed into the tree comes from a bounded source.

\subsection{From counts to a move}
\label{sec:move-selection}

After the budget is exhausted, the root's edges hold visit counts. Two outputs are
needed: a move to play, and, during self-play, a distribution to train on.

The training target is the normalized visit counts, $\pi_a = N_a / \sum_b N_b$.
Using counts rather than the mean values $Q$ is deliberate, and it is a common
implementation error: the visit count is the search's accumulated preference, and
it already embeds the value comparisons. Using $Q$ instead would discard the
search's exploration behaviour and produce a target that is closer to a greedy
policy than to the improved policy the loop requires.

The played move is drawn from $N_a^{1/\tau}$, which interpolates between uniform
sampling at large $\tau$ and argmax as $\tau \to 0$. Two implementation points
matter. First, the temperature exponent must be applied to the counts before
normalization, and the sampling must be from the resulting distribution rather
than from the raw counts, or the temperature has no effect. Second, the sampling
must be driven by an explicitly seeded generator stored per worker, so that a
self-play game is reproducible from a seed (Section~\ref{sec:determinism}).

Root noise modifies the priors, not the counts, and only at the root, and only in
self-play. Applying it at interior nodes injects randomness into every lookahead,
which degrades the search; applying it in evaluation games blurs the strength
measurement. Both errors are easy to make and hard to see, and both are caught by
the same discipline: make the noise application a function of an explicit
\code{Mode} value passed in, rather than a condition on a global flag read at the
call site.

\subsection{Reuse across moves, and the cost of re-rooting}
\label{sec:reroot}

After a move is played, the subtree below the played edge is a valid search tree
rooted at the new position: its statistics reflect real simulations, its priors
come from a valid network, and its subtree has been explored. Keeping it costs one
index assignment and gives the next decision a head start.

\begin{lstlisting}[style=rust,caption={Re-rooting. The arena is unchanged; only the
root moves.},label={lst:reroot}]
fn advance(&mut self, played: EdgeId) {
    self.root = self.edges[played as usize].child;
}
\end{lstlisting}

Three consequences must be accepted. The abandoned part of the arena -- the old
root's other subtrees -- remains allocated and unreachable; with a few thousand
nodes per search this is negligible, and it is the price of not compacting.
Second, the node's \code{parent} field is now stale for the new root and remains
valid for everything below it, so \code{parent} must not be used for anything
that assumes it reaches the root; traversal is always root-to-leaf, and the
parent link exists only for diagnostics. Third, the new root's priors carry no
noise, because noise is applied at expansion time to the node that was the root
of that search. That is correct for the next move -- noise is applied fresh to
the new root before the next search -- but it does mean that a reused tree must
not skip the noise application on the grounds that expansion already happened.

Whether to reuse at all is an experiment worth running explicitly. Reuse
improves the strength of a fixed simulation budget, and it changes the
distribution of the training targets, because a reused tree's visit counts are
the result of more simulations than the budget alone would produce. Both effects
are real; only measurement separates them.

\subsection{The evaluator boundary}
\label{sec:evaluator}

The tree needs one thing from the outside world: for a position, a policy and a
value. Making that a named trait, rather than a closure parameter, buys three
things at once.

\begin{lstlisting}[style=rust,caption={The whole coupling between search and
network. Plain data in, plain data out.},label={lst:evaluator}]
pub struct EvalRequest { pub planes: Planes, pub legal: MoveSet }
pub struct EvalResult  { pub policy: Vec<f32>, pub value: f32 }

pub trait Evaluator {
    fn evaluate(&mut self, req: EvalRequest) -> EvalResult;
}
\end{lstlisting}

First, the search becomes testable without a network. A mock implementation that
derives priors and values from the tactical predicates of
Section~\ref{sec:alphaepsilon} is deterministic and needs no framework, so the
entire search can be built and verified with no GPU present.

Second, the production implementation is genuinely stateful: it holds a channel
sender and blocks on a reply. A closure type would have to hide that state in a
captured environment, which obscures the ownership structure and makes the
blocking behaviour invisible at the call site. A named trait makes the state
explicit, and it makes the fact that \code{evaluate} may block a property of the
type rather than a comment.

Third, the search crate never depends on the deep-learning framework. The
boundary carries plain arrays and one scalar, so no tensor type, no backend type
parameter, and no device handle ever crosses into the search. This is what keeps
the search compilable and testable on a machine with no accelerator, and it is
what makes the batching service of Section~\ref{sec:parallel} possible without a
single lock on shared weights.

For a single-threaded search, a synchronous interface is the right shape even
though the implementation behind it will block on a channel. The search does not
need to know that its call crosses a thread boundary; that indirection is the
trait's entire purpose.

\subsection{What is deliberately excluded}
\label{sec:excluded}

Two classical techniques are deliberately absent, and the reasons are worth
recording because both are the obvious things to try next.

\paragraph{Transposition tables.} Two move orders can reach the same position, and
a transposition table merges them so the second arrival reuses the first's
statistics. In Gomoku the rate of transposition is low, because stones never move
once placed and the reachable positions from a fixed move order differ broadly.
More importantly, merged nodes interact badly with PUCT: the prior stored on an
edge belongs to one path, the visit counts belong to several, and the
sign convention of Section~\ref{sec:backup} assumes a unique parent. A merged
node can have several parents to whose perspective its value must be
re-presented, which is exactly the bookkeeping that the tree structure exists to
avoid. The technique is worth revisiting only if profiling shows duplicate
subtrees dominating the search.

\paragraph{Tree parallelization with virtual loss.} The classical way to
parallelize a single search tree is to let several threads descend at once and
penalize the paths under active exploration, so that followers choose elsewhere.
The penalty is a virtual loss: temporarily record a loss on the visited edges,
then undo it. AlphaGo (2016) used exactly this construction
\citep{chaslot2008parallel,segal2011paralleluct,silver2016alphago}. It is the
right answer to a problem defined by scarcity of evaluations relative to the
depth of a single search. Section~\ref{sec:parallel} argues that the situation is
inverted in a self-play system on one accelerator, and that the correct
parallelism is across games rather than within a tree. The mechanisms for the
other case are catalogued and available, but they should be introduced only when
single-game latency, not aggregate throughput, has become the binding
constraint.
